% Colocviu DCTC, proba practică, sarcina 3
% Compilați cu pdfLaTeX (implicit în Overleaf). Diacriticele se scriu direct (ă, â, î, ș, ț).
\documentclass[a4paper,11pt]{article}
\usepackage[utf8]{inputenc}
\usepackage[T1]{fontenc}
\usepackage[romanian]{babel}
\usepackage{amsmath}
\usepackage{booktabs}
\usepackage{geometry}
\geometry{margin=2.5cm}

\title{Colocviu DCTC: proba practică}
\author{Dezvoltare colaborativă și tehnologii cloud}
\date{\today}

\begin{document}
\maketitle

\section{Instrucțiuni}

Adăugați, sub această secțiune, o secțiune nouă cu titlul exact numele dumneavoastră de utilizator GitHub
(de exemplu \verb|\section{nume-utilizator}|). În ea scrieți formula variantei dumneavoastră, într-un mediu
\texttt{equation}, ca în tabelul din README:

\begin{itemize}
  \item S (simularea): $A = \pi r^2$;
  \item A: $\bar{x} = \frac{1}{n}\sum_{i=1}^{n} x_i$;
  \item B: $S = \sum_{i=1}^{n} (x_i - \bar{x})^2$;
  \item C: $x_{1,2} = \frac{-b \pm \sqrt{b^2 - 4ac}}{2a}$;
  \item D: $\sum_{k=1}^{n} k = \frac{n(n+1)}{2}$;
  \item E: $(a+b)^2 = a^2 + 2ab + b^2$;
  \item F: $n! = \prod_{k=1}^{n} k$.
\end{itemize}

Apoi adăugați în tabelul de mai jos un rând cu numele dumneavoastră de utilizator și valoarea $N$ afișată
de notebook (sarcina 4). Compilați fără erori și salvați PDF-ul ca \texttt{report/colocviu.pdf}.

% ---- Secțiunea dumneavoastră vine aici ----

\section{Tabel}

\begin{table}[h]
  \centering
  \begin{tabular}{lr}
    \toprule
    Utilizator & $N$ \\
    \midrule
    exemplu & 1234 \\
    % ---- rândul dumneavoastră vine aici ----
    \bottomrule
  \end{tabular}
  \caption{Valoarea calculată în notebook.}
\end{table}

\end{document}
